\chapter{Presentaci\'o i marc general}
\label{chap:presentacio}

\section{Presentaci\'o del curs}

Aquest curs de Programaci\'o s'adre\c{c}a a alumnat de \CourseLevel{} sense experi\`encia pr\`evia en programaci\'o i t\'e com a objectiu general introduir-lo al pensament computacional i a la programaci\'o real mitjan\c{c}ant projectes pr\`actics, progressius i contextualitzats. El curs est\`a dissenyat per desenvolupar una base s\`olida de raonament l\`ogic, descomposici\'o de problemes, expressi\'o formal d'idees i autonomia en la construcci\'o de solucions, prenent com a refer\`encia la idea del pensament computacional com a compet\`encia general per resoldre problemes \cite{wing2006} i l'aprenentatge basat en projectes com a metodologia per donar significat al treball de l'alumnat \cite{pblworks2024}.

La durada total del curs \'es de \CourseDuration{}, distribu\"ides en quatre sprints de dues setmanes cadascun, amb una c\`arrega de tres hores setmanals. Aquesta organitzaci\'o pret\'en afavorir un ritme de treball continu, exigent i assolible, orientat a la consecuci\'o d'objectius concrets.

El curs es divideix en \textbf{4 lliurables principals}, cadascun dels quals disposar\`a del seu propi document d'especificaci\'o. Per a cada lliurable, l'alumne rebr\`a un \textbf{document espec\'ific amb la definici\'o completa, l'estructura esperada i els criteris d'avaluaci\'o}, amb un m\'inim de \textbf{2 setmanes d'antelaci\'o} respecte de la data de lliurament. A mesura que el curs avanci, aquests documents s'aniran publicant i comunicant oportunament.

\section{\'Us de la Intel\textperiodcentered lig\`encia Artificial}

L'\'us d'eines d'intel\textperiodcentered lig\`encia artificial est\`a \textbf{perm\`es i \'es benvingut} en tots els lliurables del curs. L'enfocament docent parteix del sup\`osit expl\'icit que l'alumnat far\`a servir aquestes eines com a suport per explorar idees, contrastar solucions, desbloquejar dificultats i accelerar el proc\'es de treball. En conseq\"u\`encia, el \textbf{nivell d'exig\`encia} i el \textbf{ritme de treball} del curs seran superiors als que correspondrien en un context sense aquest recurs. Aquest posicionament \'es coherent amb les recomanacions actuals sobre un \'us pedag\`ogic, supervisat i responsable de la IA generativa en educaci\'o \cite{unesco2023}.

Tanmateix, l'\'us d'IA no eximeix en cap cas l'alumne de la responsabilitat acad\`emica sobre el seu treball. Tot lliurable haur\`a de poder ser compr\`es, justificat i explicat oralment o per escrit pel mateix alumne. La qualitat d'un lliurable no es valorar\`a nom\'es pel resultat final, sin\'o tamb\'e per la capacitat de demostrar comprensi\'o real sobre les decisions preses, l'estructura de la soluci\'o i el funcionament del codi o del document presentat.

\section{Pol\'itica de lliuraments}

Totes les dates de lliurament es comunicaran amb antelaci\'o en el document espec\'ific de cada activitat. \textbf{No s'acceptar\`a cap lliurament fora de termini sota cap circumst\`ancia.} Tot lliurable no presentat dins del termini establert comptar\`a com a no presentat.
